\kbvchapter{ಪಾಠ ೨೦}{ಅಲಂಕಾರಗಳು}

\indent ಬರವಣಿಗೆಯು ಅರ್ಥ ಮಾಧುರ್ಯದಿಂದ ಮತ್ತು ಶಬ್ದ ಮಾಧುರ್ಯದಿಂದ ಕೂಡಿದ್ದರೆ ಓದಲು ಕೇಳಲು ಹಿತವಾಗಿರುವುದು. ಮನಸೆಳೆಯುವ ಶಬ್ದಾರ್ಥ ಗಳಿಂದಲೂ, ಕಿವಿಗಿಂಪಾದ ಶಬ್ದಗಳಿಂದಲೂ ಬರವಣಿಗೆಯು ಕೂಡಿದ್ದರೆ ಓದ ಬೇಕೆಂಬ ಇಚ್ಛೆ ಪ್ರಬಲಿಸುವುದು. ಹೀಗೆ, ಅಲಂಕಾರಮಯವಾಗಿರುವ ಬರ ವಣಿಗೆಯು ಚಿತ್ತಾಕರ್ಷಣೀಯವಾಗಿದ್ದು ರಸದೌತಣಕ್ಕೆ ಮಿಗಿಲಾಗಿರುವುದು,\par

\indent ಶಬ್ದಾರ್ಥದಲ್ಲಿ ಚಮತ್ಕಾರವಿದ್ದು ಆ ಮೂಲಕವಾಗಿ ಬರಹದಲ್ಲಿ ಸೊಬಗು ಮೂಡಿಬಂದರೆ ಅದು ಅರ್ಥಾಲಂಕಾರದಿಂದ ಕೂಡಿದ ಬರವಣಿಗೆಯೆನಿಸುವುದು.\par

\indent ಶಬ್ದಗಳ ಪ್ರಯೋಗದಲ್ಲಿ ಜಾಣತನವಿದ್ದು ಆ ಮೂಲಕವಾಗಿ ಬರಹ ದಲ್ಲಿ ಸೊಬಗು ಮೂಡಿಬಂದರೆ, ಅದು “ಶಬ್ದಾಲಂಕಾರ”ದಿಂದ ಕೂಡಿದ ಬರವಣಿಗೆ ಎನಿಸುವುದು.\par

\vspace{7pt plus 6pt minus 1pt}\begin{center}{\kbvsecfont\fontsize{13pt}{16.5pt}\selectfont\color{kbvaccentdark} 1. ಅರ್ಥಾಲಂಕಾರ}\end{center}\vspace{3pt plus 3pt minus 1pt}

\indent ಅರ್ಥಾಲಂಕಾರದಲ್ಲಿ ಹಲವು ಪ್ರಕಾರಗಳುಂಟು. ಅವುಗಳು ಯಾವು ವೆಂಬುದನ್ನು ನೋಡೋಣ.\par

\kbvsechead{(೧)\hspace{0.45em}ಉಪಮಾಲಂಕಾರ :}

\indent “ಇಂದ್ರಾಣಿಯ ಮುಖವು ತಾವರೆಯ ಹೂವಿನಂತೆ ಕಂಗೊಳಿಸುತ್ತಿತ್ತು.”\par

\noindent\hfill{\fontsize{9.5pt}{12.5pt}\selectfont\color{kbvgrey}- ಈ ವಾಕ್ಯದಲ್ಲಿ ಇಂದ್ರಾಣಿಯ ಮುಖವನ್ನು ತಾವರೆಯ ಹೂವಿಗೆ ಹೋಲಿಸಿ}\par\vspace{3pt plus 2pt}

\indent ರುವುದನ್ನು ಕಾಣುತ್ತೇವೆ. ಇಂತಹ ಹೋಲಿಕೆಗಳಲ್ಲಿ ಯಾವುದಕ್ಕೆ ಹೋಲಿಕೆ\par

\indent ಯನ್ನು ಕೊಡಲಾಗುವುದೋ ಅದನ್ನು “ಉಪಮೇಯ” ಎಂದು, ಯಾವು\par

\indent ದನ್ನು ಹೋಲಿಸಲಾಗುವುದೋ ಅದನ್ನು “ಉಪಮಾನ” ಎಂದು ಕರೆಯು ವುದು ವಾಡಿಕೆಯಲ್ಲಿದೆ. ಈ ಉದಾಹರಣೆಯಲ್ಲಿ “ಇಂದ್ರಾಣಿಯ ಮುಖ” ಉಪಮೇಯವಾಗಿರುವುದು, “ತಾವರೆಯ ಹೂವು” ಉಪಮಾನವಾಗಿರುವುದು. ಈ ಬಗೆಯ ಅರ್ಥಾಲಂಕಾರವನ್ನು ಉಪಮಾಲಂಕಾರವೆನ್ನುವರು.\par
